\documentclass[10pt,landscape]{article}
\usepackage[landscape,margin=0.5in]{geometry}
\usepackage{multicol}
\usepackage{amsmath,amssymb,amsfonts}
\usepackage{enumitem}
\usepackage{titlesec}
\usepackage{fancyhdr}
\usepackage{xcolor}

\pagestyle{empty}
\setlength{\parindent}{0pt}
\setlength{\parskip}{2pt}

\titleformat{\section}{\bfseries\normalsize}{}{0em}{}[\titlerule]
\titleformat{\subsection}{\bfseries\small}{}{0em}{}
\titlespacing{\section}{0pt}{4pt}{2pt}
\titlespacing{\subsection}{0pt}{3pt}{1pt}

\begin{document}
\begin{multicols}{3}

\begin{center}
{\large\bfseries CSCE 763 Cheat Sheet}
\end{center}

\section{Image Operations}
\small
Example image: \textit{firework burst against a dark night sky.}

\subsection{Intensity Transformations ($s = T(r)$)}
\textbf{Negative / Complement} ($s = L{-}1{-}r$):
Inverts all intensities. The dark sky becomes bright white; the bright firework sparks become dark. Useful for enhancing detail in predominantly dark images.

\textbf{Log Transform} ($s = c\log(1+r)$):
Compresses high-intensity values and expands low ones. Bright firework core compresses; faint glow in the sky becomes more visible. Good for displaying wide dynamic range.

\textbf{Power-Law / Gamma} ($s = cr^\gamma$):
$\gamma < 1$: brightens dark regions---faint smoke and sky glow become visible.
$\gamma > 1$: darkens the image---only the brightest sparks remain visible, sky becomes pure black.

\textbf{Contrast Stretching}:
Maps a narrow intensity range to the full $[0, L{-}1]$ range. Increases separation between the dim sky and bright sparks, making the firework pop.

\textbf{Intensity-Level Slicing}:
Highlights a specific intensity band. Can isolate just the bright firework trails (high intensities) while suppressing or preserving the dark background.

\textbf{Bit-Plane Slicing}:
Decomposes into 8 binary planes. Higher bit planes (7, 8) capture the firework structure; lower planes (1, 2) are mostly noise from the dark sky.

\textbf{Thresholding} ($s = 0$ or $L{-}1$):
Produces binary image. Separates bright firework pixels from the dark sky into pure white/black.

\subsection{Histogram Operations}
\textbf{Histogram Equalization}:
Flattens the histogram to use full intensity range. The firework image (mostly dark pixels) gets its few bright values spread out, revealing more contrast in the sparks and smoke trails.

\textbf{Histogram Matching}:
Reshapes histogram to match a target distribution. Can make the firework image's tonal distribution match a reference image.

\textbf{Local Histogram Processing}:
Applies equalization in sliding neighborhoods. Reveals faint details in both the dark sky regions and around the bright firework center that global methods miss.

\subsection{Arithmetic Operations}
\textbf{Addition / Averaging}:
Average $K$ noisy firework frames: $\bar{g}(x,y) = \frac{1}{K}\sum g_i(x,y)$. Reduces noise variance by $1/K$; the firework sharpens while random noise fades.

\textbf{Subtraction}:
Subtract a sky-only frame from the firework frame. Isolates just the firework burst by removing the background. Images must be registered.

\textbf{Multiplication (Masking)}:
$g = f \cdot h$: multiply by an ROI mask to isolate only the firework region. Division corrects uneven illumination / shading.

\subsection{Spatial Filtering}
\textbf{Smoothing (Box / Gaussian)}:
Averages neighborhood pixels. Firework edges soften; sparks blur into the sky. Reduces noise but loses fine spark detail.

\textbf{Median Filter} (nonlinear):
Replaces each pixel with the neighborhood median. Removes salt-and-pepper noise from the dark sky while preserving the sharp edges of firework trails better than averaging.

\textbf{Sharpening (Laplacian, Unsharp Mask)}:
Enhances edges and fine detail. Firework spark tips and trail edges become crisper and more defined against the sky.

\subsection{Geometric Operations}
\textbf{Scaling / Resizing}:
Enlarging requires interpolation---NN gives blocky sparks, bilinear smooths them, bicubic preserves fine trails best.

\textbf{Rotation / Translation}:
Repositions the firework in the frame. Requires resampling via interpolation at non-integer coordinates.

\subsection{Logic Operations (Binary)}
\textbf{NOT}: swaps foreground/background.
\textbf{AND}: intersection of two masks---e.g., where firework overlaps a region of interest.
\textbf{OR}: union of two binary firework masks.
\textbf{XOR}: pixels in one mask but not both.

\end{multicols}
\end{document}
